\documentclass[main.tex]{subfiles}


\begin{document}

%from pagenumber to up
% \phantomsection 
% \hypertarget{toc}{}

 

 

% номер секции вручную
\setcounter{section}{47
}
\addtocounter{section}{-1} 

\section{Теплопередача}


\textbf{Теплопередача}~--- это процесс перехода внутренней энергии от более горячего тела к более холодному без совершения работы. Этот процесс также называют \emph{теплообменом}. Например, если в холодный стакан налить горячую воду, то стакан нагреется~--- произошла теплопередача от воды к стакану.   

Различают три вида теплопередачи.

\sbox{0}{%
    \begin{tikzpicture}[font=\small,            line cap=round,                             >={Stealth[inset=0pt,round,length=3mm, angle'=20]},vec/.style={>={Stealth[inset=0pt,round,length=3.5mm, angle'=20]}}]


% plane
\filldraw[lightgray,semitransparent] (0,0) -- (5,0) --++ (0,-.3) --++ (-5,0) -- cycle;
\draw[] (0,0) -- (5,0);


%solid
\fill[lightgray] (2.5,1) rectangle++ (2,.2);
    % red part
    \fill[red,path fading=east] (2.5,1) rectangle++ (1,.2);
    \draw[] (2.5,1) rectangle++ (2,.2);


%flame

\begin{scope}[yshift=.1cm]
    
\fill[red,nearly opaque] (2.5,0) 
to [out=180,in=-90] (2.2,.3) 
to [out=90,in=-90] (2.5,1) 
to [out=-45,in=90] (2.7,.5) 
to [out=45,in=-90] (2.75,.7) 
to [out=-45,in=0] (2.5,0);

\fill[scale around={.8:(2.5,0)},orange] (2.5,0) 
to [out=180,in=-90] (2.2,.3) 
to [out=90,in=-90] (2.55,1) 
to [out=-45,in=90] (2.65,.4) 
to [out=45,in=-90] (2.8,.55) 
to [out=-45,in=0] (2.5,0);

\fill[scale around={.5:(2.5,0)},yellow] (2.5,0) 
to [out=180,in=-90] (2.2,.3) 
to [out=90,in=-110] (2.55,1) 
to [out=-70,in=150] (2.8,.5) 
to [out=-100,in=90] (2.8,0.3)
to [out=-90,in=0] (2.5,0);

\end{scope}

\filldraw[] (2,0) rectangle++ (1,.1);




    \end{tikzpicture}%
}

\nointerlineskip
\task{%
{%
\setlength\intextsep{0pt}
\begin{wrapfigure}[6]{r}{\wd0}
    \centering
    \usebox{0}%
    \caption{Теплопроводность}
    \label{fig:p77}
\end{wrapfigure}
\textbf{Теплопроводность}~--- это теплообмен между контактирующими телами (или частями тела). На рис.\,\ref{fig:p77} изображен металлический стержень, помещенный одним концом в огонь.
    
%         \begin{figure}[H]
%     \centering

%     \begin{tikzpicture}[font=\small,            line cap=round,                             >={Stealth[inset=0pt,round,length=3mm, angle'=20]},vec/.style={>={Stealth[inset=0pt,round,length=3.5mm, angle'=20]}}]


% % plane
% \filldraw[lightgray,semitransparent] (0,0) -- (5,0) --++ (0,-.3) --++ (-5,0) -- cycle;
% \draw[] (0,0) -- (5,0);


% %solid
% \fill[lightgray] (2.5,1) rectangle++ (2,.2);
%     % red part
%     \fill[red,path fading=east] (2.5,1) rectangle++ (1,.2);
%     \draw[] (2.5,1) rectangle++ (2,.2);


% %flame

% \begin{scope}[yshift=.1cm]
    
% \fill[red,nearly opaque] (2.5,0) 
% to [out=180,in=-90] (2.2,.3) 
% to [out=90,in=-90] (2.5,1) 
% to [out=-45,in=90] (2.7,.5) 
% to [out=45,in=-90] (2.75,.7) 
% to [out=-45,in=0] (2.5,0);

% \fill[scale around={.8:(2.5,0)},orange] (2.5,0) 
% to [out=180,in=-90] (2.2,.3) 
% to [out=90,in=-90] (2.55,1) 
% to [out=-45,in=90] (2.65,.4) 
% to [out=45,in=-90] (2.8,.55) 
% to [out=-45,in=0] (2.5,0);

% \fill[scale around={.5:(2.5,0)},yellow] (2.5,0) 
% to [out=180,in=-90] (2.2,.3) 
% to [out=90,in=-110] (2.55,1) 
% to [out=-70,in=150] (2.8,.5) 
% to [out=-100,in=90] (2.8,0.3)
% to [out=-90,in=0] (2.5,0);

% \end{scope}

% \filldraw[] (2,0) rectangle++ (1,.1);




%     \end{tikzpicture}
% \caption{Теплопроводность}
% \label{fig:p77}
% \end{figure}

    
    % Ясно, что температура соседних частей стержня (рис.\,\ref{fig:p77}) будет постепенно повышаться. 

    \smallskip
    
    <<Соприкасающиеся>> с огнем частицы стержня  (молекулы или атомы) начинают интенсивнее колебаться и сильнее <<толкают>> соседние частицы. Тепловое движение соседних частиц также увеличивается, и они в свою очередь <<раскачивают>> уже своих соседей по другую сторону. Так тепло постепенно распространяется от участка к участку.

}

}

%%%%%%%%%%%%%%%%%%%%%%%%%%%%%%%%%%%

\tikzfading[name=fade north,
top color=transparent!100,
bottom color=transparent!70]

\sbox{0}{%
        \begin{tikzpicture}[font=\small,            line cap=round,                             >={Stealth[inset=0pt,round,length=3mm, angle'=20]},vec/.style={>={Stealth[inset=0pt,round,length=3.5mm, angle'=20]}}]


% plane
\filldraw[lightgray,semitransparent] (0,0) -- (5,0) --++ (0,-.3) --++ (-5,0) -- cycle;
\draw[] (0,0) -- (5,0);


% %solid
% \fill[lightgray] (2.5,1) rectangle++ (2,.2);
%     % red part
%     \fill[red,path fading=east] (2.5,1) rectangle++ (1,.2);
%     \draw[] (2.5,1) rectangle++ (2,.2);


%water
\draw (1.5,1.1) coordinate (w1) ++ (0,1.75) --++ (0,-1.75) --++ (2,0) --++ (0,1.75);

\fill [SkyBlue,semitransparent] (w1) rectangle++ (2,1.5);
\fill[red,path fading=fade north] (w1) rectangle++ (2,.75);

\draw[->,semithick,vec,red] (2.9,1.2) 
to [out=180,in=-120] (2.7,2.2);
\draw[->,semithick,vec,NavyBlue,shorten <=.05cm,shorten >=.05cm] (2.7,2.2) 
to [out=60,in=180] (3.05,2.45)
to [out=0,in=0] (2.9,1.2);

\draw[->,semithick,vec,red] (2.1,1.2) 
to [out=0,in=-60] (2.3,2.2);
\draw[->,semithick,vec,NavyBlue,shorten <=.05cm,shorten >=.05cm] (2.3,2.2) 
to [out=120,in=0] (1.95,2.45)
to [out=180,in=180] (2.1,1.2);


%flame

\begin{scope}[yshift=.1cm]
    
\fill[red,nearly opaque] (2.5,0) 
to [out=180,in=-90] (2.2,.3) 
to [out=90,in=-90] (2.5,1) 
to [out=-45,in=90] (2.7,.5) 
to [out=45,in=-90] (2.75,.7) 
to [out=-45,in=0] (2.5,0);

\fill[scale around={.8:(2.5,0)},orange] (2.5,0) 
to [out=180,in=-90] (2.2,.3) 
to [out=90,in=-90] (2.55,1) 
to [out=-45,in=90] (2.65,.4) 
to [out=45,in=-90] (2.8,.55) 
to [out=-45,in=0] (2.5,0);

\fill[scale around={.5:(2.5,0)},yellow] (2.5,0) 
to [out=180,in=-90] (2.2,.3) 
to [out=90,in=-110] (2.55,1) 
to [out=-70,in=150] (2.8,.5) 
to [out=-100,in=90] (2.8,0.3)
to [out=-90,in=0] (2.5,0);

\end{scope}

\filldraw[] (2,0) rectangle++ (1,.1);




    \end{tikzpicture}%
}

\nointerlineskip
\task{%
{%
\setlength\intextsep{0pt}
\begin{wrapfigure}[9]{r}{\wd0}
    \centering
    \usebox{0}%
    \caption{Конвекция}
    \label{fig:p78}
\end{wrapfigure}
\textbf{Конвекция}~--- это теплообмен в жидкостях или газах за счет потоков вещества. На рис.\,\ref{fig:p78} изображен сосуд с водой над огнем.

% \tikzfading[name=fade north,
% top color=transparent!100,
% bottom color=transparent!70]

%             \begin{figure}[H]
%     \centering

%     \begin{tikzpicture}[font=\small,            line cap=round,                             >={Stealth[inset=0pt,round,length=3mm, angle'=20]},vec/.style={>={Stealth[inset=0pt,round,length=3.5mm, angle'=20]}}]


% % plane
% \filldraw[lightgray,semitransparent] (0,0) -- (5,0) --++ (0,-.3) --++ (-5,0) -- cycle;
% \draw[] (0,0) -- (5,0);


% % %solid
% % \fill[lightgray] (2.5,1) rectangle++ (2,.2);
% %     % red part
% %     \fill[red,path fading=east] (2.5,1) rectangle++ (1,.2);
% %     \draw[] (2.5,1) rectangle++ (2,.2);


% %water
% \draw (1.5,1.1) coordinate (w1) ++ (0,2) --++ (0,-2) --++ (2,0) --++ (0,2);

% \fill [SkyBlue,semitransparent] (w1) rectangle++ (2,1.5);
% \fill[red,path fading=fade north] (w1) rectangle++ (2,.75);

% \draw[->,semithick,vec,red] (2.9,1.2) 
% to [out=180,in=-120] (2.7,2.2);
% \draw[->,semithick,vec,NavyBlue,shorten <=.05cm,shorten >=.05cm] (2.7,2.2) 
% to [out=60,in=180] (3.05,2.45)
% to [out=0,in=0] (2.9,1.2);

% \draw[->,semithick,vec,red] (2.1,1.2) 
% to [out=0,in=-60] (2.3,2.2);
% \draw[->,semithick,vec,NavyBlue,shorten <=.05cm,shorten >=.05cm] (2.3,2.2) 
% to [out=120,in=0] (1.95,2.45)
% to [out=180,in=180] (2.1,1.2);


% %flame

% \begin{scope}[yshift=.1cm]
    
% \fill[red,nearly opaque] (2.5,0) 
% to [out=180,in=-90] (2.2,.3) 
% to [out=90,in=-90] (2.5,1) 
% to [out=-45,in=90] (2.7,.5) 
% to [out=45,in=-90] (2.75,.7) 
% to [out=-45,in=0] (2.5,0);

% \fill[scale around={.8:(2.5,0)},orange] (2.5,0) 
% to [out=180,in=-90] (2.2,.3) 
% to [out=90,in=-90] (2.55,1) 
% to [out=-45,in=90] (2.65,.4) 
% to [out=45,in=-90] (2.8,.55) 
% to [out=-45,in=0] (2.5,0);

% \fill[scale around={.5:(2.5,0)},yellow] (2.5,0) 
% to [out=180,in=-90] (2.2,.3) 
% to [out=90,in=-110] (2.55,1) 
% to [out=-70,in=150] (2.8,.5) 
% to [out=-100,in=90] (2.8,0.3)
% to [out=-90,in=0] (2.5,0);

% \end{scope}

% \filldraw[] (2,0) rectangle++ (1,.1);




%     \end{tikzpicture}
% \caption{Конвекция}
% \label{fig:p78}
% \end{figure}

    \smallskip
    
    Жидкость вблизи дна сосуда нагревается и расширяется, так что ее плотность становится меньше по сравнению с холодной жидкостью сверху. Эта менее плотная теплая жидкость под действием силы Архимеда
    % \footnote{Точнее~--- под действием сил тяжести и Архимеда.} 
    поднимается вверх (красные стрелки на рис.\,\ref{fig:p78}); напротив, более холодная жидкость <<тонет>>, то есть опускается вниз (синие стрелки на рис.\,\ref{fig:p78}). Таким образом вещество в сосуде перемешивается, и вся жидкость с течением времени прогревается.  
   
}

}

%%%%%%%%%%%%%%%%%%%%%%%%%%%%%%%%

\tikzfading[name=fade out,
inner color=transparent!0,
outer color=transparent!100]

\sbox{0}{%
        \begin{tikzpicture}[font=\small,            line cap=round,                             >={Stealth[inset=0pt,round,length=3mm, angle'=20]},vec/.style={>={Stealth[inset=0pt,round,length=3.5mm, angle'=20]}}]


% plane
\filldraw[lightgray,semitransparent] (0,0) -- (5,0) --++ (0,-.3) --++ (-5,0) -- cycle;
\draw[] (0,0) -- (5,0);


% %solid
% \fill[lightgray] (2.5,1) rectangle++ (2,.2);
%     % red part
%     \fill[red,path fading=east] (2.5,1) rectangle++ (1,.2);
%     \draw[] (2.5,1) rectangle++ (2,.2);

%radiation
\begin{scope}
    \clip (4.5,0) arc [start angle=0, end angle=180, radius=2] -- cycle;
    \fill [yellow,path fading=fade out] (.5,-2) rectangle++ (4,4);
\end{scope}


%flame

\begin{scope}[yshift=.1cm]
    
\fill[red,nearly opaque] (2.5,0) 
to [out=180,in=-90] (2.2,.3) 
to [out=90,in=-90] (2.5,1) 
to [out=-45,in=90] (2.7,.5) 
to [out=45,in=-90] (2.75,.7) 
to [out=-45,in=0] (2.5,0);

\fill[scale around={.8:(2.5,0)},orange] (2.5,0) 
to [out=180,in=-90] (2.2,.3) 
to [out=90,in=-90] (2.55,1) 
to [out=-45,in=90] (2.65,.4) 
to [out=45,in=-90] (2.8,.55) 
to [out=-45,in=0] (2.5,0);

\fill[scale around={.5:(2.5,0)},yellow] (2.5,0) 
to [out=180,in=-90] (2.2,.3) 
to [out=90,in=-110] (2.55,1) 
to [out=-70,in=150] (2.8,.5) 
to [out=-100,in=90] (2.8,0.3)
to [out=-90,in=0] (2.5,0);

\end{scope}

\filldraw[] (2,0) rectangle++ (1,.1);




    \end{tikzpicture}%
}

\nointerlineskip
\task{%
{%
\setlength\intextsep{0pt}
\begin{wrapfigure}[]{r}{\wd0}
    \centering
    \usebox{0}%
    \caption{Излучение}
    \label{fig:p79}
\end{wrapfigure}
\textbf{Излучение}~--- это теплообмен посредством электромагнитных волн. На рис.\,\ref{fig:p79}~--- огонь.

% \tikzfading[name=fade north,
% top color=transparent!100,
% bottom color=transparent!70]

%             \begin{figure}[H]
%     \centering

%     \begin{tikzpicture}[font=\small,            line cap=round,                             >={Stealth[inset=0pt,round,length=3mm, angle'=20]},vec/.style={>={Stealth[inset=0pt,round,length=3.5mm, angle'=20]}}]


% % plane
% \filldraw[lightgray,semitransparent] (0,0) -- (5,0) --++ (0,-.3) --++ (-5,0) -- cycle;
% \draw[] (0,0) -- (5,0);


% % %solid
% % \fill[lightgray] (2.5,1) rectangle++ (2,.2);
% %     % red part
% %     \fill[red,path fading=east] (2.5,1) rectangle++ (1,.2);
% %     \draw[] (2.5,1) rectangle++ (2,.2);


% %water
% \draw (1.5,1.1) coordinate (w1) ++ (0,2) --++ (0,-2) --++ (2,0) --++ (0,2);

% \fill [SkyBlue,semitransparent] (w1) rectangle++ (2,1.5);
% \fill[red,path fading=fade north] (w1) rectangle++ (2,.75);

% \draw[->,semithick,vec,red] (2.9,1.2) 
% to [out=180,in=-120] (2.7,2.2);
% \draw[->,semithick,vec,NavyBlue,shorten <=.05cm,shorten >=.05cm] (2.7,2.2) 
% to [out=60,in=180] (3.05,2.45)
% to [out=0,in=0] (2.9,1.2);

% \draw[->,semithick,vec,red] (2.1,1.2) 
% to [out=0,in=-60] (2.3,2.2);
% \draw[->,semithick,vec,NavyBlue,shorten <=.05cm,shorten >=.05cm] (2.3,2.2) 
% to [out=120,in=0] (1.95,2.45)
% to [out=180,in=180] (2.1,1.2);


% %flame

% \begin{scope}[yshift=.1cm]
    
% \fill[red,nearly opaque] (2.5,0) 
% to [out=180,in=-90] (2.2,.3) 
% to [out=90,in=-90] (2.5,1) 
% to [out=-45,in=90] (2.7,.5) 
% to [out=45,in=-90] (2.75,.7) 
% to [out=-45,in=0] (2.5,0);

% \fill[scale around={.8:(2.5,0)},orange] (2.5,0) 
% to [out=180,in=-90] (2.2,.3) 
% to [out=90,in=-90] (2.55,1) 
% to [out=-45,in=90] (2.65,.4) 
% to [out=45,in=-90] (2.8,.55) 
% to [out=-45,in=0] (2.5,0);

% \fill[scale around={.5:(2.5,0)},yellow] (2.5,0) 
% to [out=180,in=-90] (2.2,.3) 
% to [out=90,in=-110] (2.55,1) 
% to [out=-70,in=150] (2.8,.5) 
% to [out=-100,in=90] (2.8,0.3)
% to [out=-90,in=0] (2.5,0);

% \end{scope}

% \filldraw[] (2,0) rectangle++ (1,.1);




%     \end{tikzpicture}
% \caption{Конвекция}
% \label{fig:p78}
% \end{figure}

    \smallskip
    
    Атомы любого тела при ненулевой температуре (в кельвинах) вследствие теплового движения испускают так называемые \emph{электромагнитные волны}\footnotemark, то есть излучают энергию. \emph{Видимый свет}~--- это частный случай излучения (свечение огня на рис.\,\ref{fig:p79} является примером видимого излучения). Тела излучают энергию во все стороны; однако, не всякое излучение можно видеть~--- ведь большинство предметов при нормальных условиях не светятся, хотя эти предметы несомненно излучают (пример~--- горячая печь). Именно излучение <<доставляет>> энергию от Солнца к планетам~--- излучение может распространяться как в веществе, так и в \emph{вакууме}.
   
}

}

% \par\addvspace{\bigskipamount} % отступ подобно нумерованным окружениям

% \footnotetext{Точнее говоря, электромагнитные волны излучают \emph{заряженные частицы}, входящие в состав атомов и совершающие вместе с ними хаотическое движение. Эти волны называют еще \emph{тепловым излучением}~--- в напоминание о том, что их источником служит тепловое движение частиц вещества.}

\footnotetext{Электромагнитная волна~--- это распространение колебаний \emph{электрического} и \emph{магнитного полей}~--- особых форм материи, окружающей движущиеся \emph{заряженные частицы}. Именно такие частицы, входящие в состав атома и совершающие вместе с ними хаотическое движение, и излучают электромагнитные волны, уносящие с собой часть внутренней энергии тела.}




















 
\end{document}